\chapter{예측 모델과 도메인 지식 처치}
\label{ch:models}

모든 모델은 동일한 관측 창 텐서를 입력으로 받아 로짓 $z$를 출력하고 $p = \sigma(z)$를 블루 팀 승리 확률로 해석한다. 모델은 입력 \emph{표현}에 따라 테이블, 시퀀스, 그래프, 시공간 그래프, 융합의 다섯 계열로 나뉜다. 계열별 대표 하이퍼파라미터는 부록 \ref{app:hparams}에 있다.

\section{테이블 기준선: LightGBM}
\label{sec:model-lgbm}

식 \eqref{eq:phitab}의 테이블 벡터에 LightGBM \cite{ke2017lightgbm}을 적용한다. 5{,}000 라운드, 학습률 0.03, 최대 깊이 6, 잎 31, 행·열 부표본 0.7, $L_1$ 1.0, $L_2$ 5.0, 잎당 최소 200 표본, 검증 AUC 기준 200 라운드 조기 종료를 사용한다. 패치 공변량 이동에 대응하여 표본 $i$의 패치 인덱스 $p_i$에 따라
\begin{equation}
  w_i = \exp\!\big((p_i - p_{\min})/\tau\big), \qquad \tau = 2.0
\end{equation}
의 최근성 가중치를 부여한다. 이는 중요도 가중 경험적 위험 최소화 \cite{shimodaira2000improving}의 단순화된 형태이다.

\section{정합 입력 다층 퍼셉트론}
\label{sec:model-mlp}

RQ3(학습기 효과)을 분리하기 위해 LightGBM과 \emph{같은} 약 2{,}980 차원 벡터를 입력으로 하는 다층 퍼셉트론(MLP)을 둔다. 이 모델은 표현을 고정한 채 학습기만 바꾼 것이므로, 두 모델의 차이는 오직 귀납 편향(트리 앙상블 대 매끄러운 함수 근사)에서 온다. 심층 하네스 안의 진단용 MLP(95 차원 거시 시퀀스의 $[x_L \,\|\, \mathrm{mean}]$ 요약 사용)는 입력이 다르므로 이 비교에 쓰지 않는다.

\section{순차 모델}
\label{sec:model-seq}

거시 시퀀스 $S \in \R^{L \times D}$($L=6$, 입력 투영 후 $D = 256$)를 처리한다.

\begin{itemize}[leftmargin=1.5em]
  \item \textbf{BiGRU / BiLSTM.} 2 층 양방향, 은닉 128, 드롭아웃 0.20. GRU 셀은 갱신 게이트 $z_t$, 재설정 게이트 $r_t$로 $h_t = (1-z_t)\odot h_{t-1} + z_t \odot \tilde h_t$를 계산하고 \cite{cho2014learning}, LSTM은 망각·입력·출력 게이트와 셀 상태를 갖는다 \cite{hochreiter1997long}. 마지막 은닉 상태(또는 T3의 주의 풀링)를 분류 헤드에 넣는다.
  \item \textbf{Transformer.} $d_{\mathrm{model}} = 64$, 헤드 4, 2 층, 피드포워드 128, 사인 위치 부호화. $\mathrm{Attention}(Q,K,V) = \softmax(QK^\top/\sqrt{d_k})V$ \cite{vaswani2017attention}.
  \item \textbf{TCN.} 채널 64, 커널 3, 팽창 $[1,2,4]$의 3 단 인과 합성곱 블록 \cite{bai2018empirical}. 수용 영역 $2(k-1)\sum d + 1 = 29 > L$.
  \item \textbf{Mamba.} $d_{\mathrm{model}} = 128$, 3 층, $d_{\mathrm{state}} = 16$, $d_{\mathrm{conv}} = 4$의 선택적 상태공간 모델 \cite{gu2023mamba}. $h_t = \bar A_t h_{t-1} + \bar B_t x_t$, $y_t = C h_t + D x_t$이며 $\bar A_t, \bar B_t$는 입력 의존적이다.
  \item \textbf{하이브리드 $h_0$ 조건화.} 테이블 요약 $\phi_{\mathrm{tab}}(S)$를 MLP로 투영하여 순환망의 초기 은닉 상태로 주입한다($h_0 = \mathrm{MLP}(\phi_{\mathrm{tab}}(S))$).
\end{itemize}

\section{그래프 신경망}
\label{sec:model-gnn}

식 \eqref{eq:adjacency}의 그래프 위에서 노드 표현을 갱신하고, 팀별 풀링 $[h_{\mathrm{blue}} \,\|\, h_{\mathrm{red}}]$ 후 분류한다. 정점 차원 96, 드롭아웃 0.25, LayerNorm.

\begin{itemize}[leftmargin=1.5em]
  \item \textbf{GCN} \cite{kipf2017semi}: $H' = \mathrm{LN}\big(\sigma(\tilde D^{-1/2}\tilde A\tilde D^{-1/2} H W) + H\big)$, $\tilde A = A + I$.
  \item \textbf{GraphSAGE} \cite{hamilton2017inductive}: $h_i' = W[h_i \,\|\, \mathrm{mean}_{j \in \mathcal{N}(i)} h_j]$.
  \item \textbf{GATv2} \cite{brody2022how}: $\alpha_{ij} = \softmax_j\big(a^\top \mathrm{LeakyReLU}(W[h_i \,\|\, h_j])\big)$, 4 헤드, $A_{ij} = 0$인 쌍은 주의에서 강제로 제외(hard mask).
  \item \textbf{MPNN} \cite{gilmer2017neural}: 변 조건부 메시지 $m_{ij} = \mathrm{MLP}_e(e_{ij}) \odot W h_j$.
  \item \textbf{GraphTransformer}: 변 가중치를 편향으로 갖는 다중 헤드 자기 주의 층.
\end{itemize}

\section{시공간 모델}
\label{sec:model-st}

\begin{itemize}[leftmargin=1.5em]
  \item \textbf{ST-GNN}: 구간마다 GNN을 적용한 뒤 시간 축을 GRU 또는 주의 풀링으로 요약.
  \item \textbf{ST-GCN} \cite{yan2018spatial}: 그래프 합성곱과 시간 합성곱 블록을 교대로 쌓음.
  \item \textbf{ST-Mamba}: GNN 공간 인코딩 후 Mamba 시간 인코딩.
  \item \textbf{다중 스케일 동적 그래프}: $A = \sum_k w_k A(\sigma_k)$의 학습 가중 커널 합.
  \item \textbf{EventXAttn}: 사건 토큰 임베딩 $e_i = \mathrm{Embed}(\mathrm{type}_i) + \mathrm{Linear}(\mathrm{cont}_i)$를 질의로, 그래프 인코딩된 플레이어 상태를 키·값으로 하는 교차 주의 후 중요도 가중 풀링.
\end{itemize}

\section{계층적 융합}
\label{sec:model-fusion}

계층적 융합(layered fusion)은 세 스트림, 즉 전역 인코더($S \to h_{\mathrm{glob}}$; GRU·LSTM·Transformer·TCN·Mamba 중 택일), 그래프 인코더($X^{\mathrm{node}}, A \to h_{\mathrm{graph}}$; GCN·GraphSAGE·GraphTransformer·GATv2·MPNN 중 택일), 사건 인코더($E \to h_{\mathrm{ev}}$; 자기 주의·교차 주의·평균 풀링 중 택일)를 게이트로 결합한다:
\begin{equation}
  h_{\mathrm{cat}} = [h_{\mathrm{glob}} \,\|\, h_{\mathrm{graph}} \,\|\, h_{\mathrm{ev}}], \qquad
  g = \sigma(W_g h_{\mathrm{cat}} + b_g), \qquad
  h_{\mathrm{fuse}} = g \odot \mathrm{MLP}(h_{\mathrm{cat}}).
  \label{eq:fusion}
\end{equation}
융합 차원 192, 게이트 은닉 64, 사건 인코더 차원 128. 선택적으로 LightGBM 로짓을 이어 붙여 $h_{\mathrm{final}} = [h_{\mathrm{fuse}} \,\|\, z_{\mathrm{lgbm}}]$로 스태킹할 수 있다(Layered + Logit). 식 \eqref{eq:fusion}의 게이트는 전문가 혼합 \cite{shazeer2017outrageously,ma2018modeling}의 소프트 게이트를 특징 차원별로 적용한 것이다.

\section{스태킹 앙상블과 사후 보정}
\label{sec:model-stacking}

기저 모델 로짓을 입력으로 하는 로지스틱 회귀 메타 학습기 \cite{wolpert1992stacked}를 (i) 학습 분할 로짓으로, (ii) 5-겹 교차검증의 out-of-fold 로짓으로, (iii) $M$개 기저 모델의 모든 비공집합 $2^M - 1$개를 열거하여 검증 AUC가 최대인 부분집합으로 적합하는 세 방식을 지원한다. 선택 후 메타 학습기를 학습+검증 분할에 재적합하고, 필요 시 패치별 온도 $T_p^* = \argmin_T \sum_i \mathcal{L}_{\mathrm{BCE}}(\sigma(z_i/T), y_i)$로 사후 보정 \cite{guo2017calibration}한다.

\section{도메인 지식 처치 T1--T7}
\label{sec:model-treatments}

RQ5의 일곱 처치는 각각 LoL 도메인 지식과 학습 이론의 근거를 갖는 독립적인 설계 요소이며, 기본 구성에서 하나씩 켜서 효과를 측정한다.

\begin{description}[leftmargin=2.5em,style=nextline]
  \item[T1 초점 손실 \cite{lin2017focal}] $\mathcal{L}_{\mathrm{FL}} = -\alpha_t (1-p_t)^\gamma \log p_t$, $\gamma = 2$, $\alpha = 0.25$. 골드 격차가 큰 ``압도'' 교전은 쉬운 표본이므로 그 기여를 줄여 대등한 교전에 학습을 집중시킨다. $p_t = 0.9$인 표본의 가중치는 $0.01$이다.
  \item[T2 게임 국면 부호화] $\phi(t) = [\sigma((14-t)/\tau),\ \sigma((t-10)/\tau)\,\sigma((28-t)/\tau),\ \sigma((t-22)/\tau)]$, $\tau = 3$ 분. 초반(라인전)·중반(오브젝트 경합)·후반(5:5 한타)을 매끄러운 3 차원 벡터로 전역 특징에 추가한다.
  \item[T3 시간 주의 풀링 \cite{bahdanau2015neural}] $e_t = w^\top\tanh(W_a h_t + b_a)$, $\alpha_t = \softmax(e)_t$, $c = \sum_t \alpha_t h_t$, 출력 $[h_L \,\|\, c]$. 마지막 은닉 상태만 쓰는 대신 파워 스파이크 등 중요한 시점을 선택적으로 참조한다.
  \item[T4 모멘텀 특징] 금융 기술 분석의 MACD \cite{appel2005technical}에서 착안하여 1차 차분 $\Delta x_t$의 단기 평균($k = 3$)과 장기 평균의 차 $\mu_{\mathrm{short}} - \mu_{\mathrm{long}}$을 테이블 특징에 추가한다. 추세와 가속을 구분한다.
  \item[T5 역할 인식 인접행렬] $A'_{ij} = A_{ij}\,\softplus(R_{\mathrm{role}(i),\mathrm{role}(j)})$, $R \in \R^{5\times 5}$ 학습, 초기값 0($\softplus(0) = \ln 2$). 같은 거리라도 정글--미드 관계가 탑--서포터 관계보다 중요할 수 있음을 학습하게 한다.
  \item[T6 다중 과제 손실 \cite{caruana1997multitask,ruder2017overview}] $\mathcal{L} = \mathcal{L}_{\mathrm{fight}} + \lambda_g\|\hat y_{\mathrm{gold}} - y_{\mathrm{gold}}\|^2 + \lambda_k\|\hat y_{\mathrm{kill}} - y_{\mathrm{kill}}\|^2 + \lambda_o\|\hat y_{\mathrm{obj}} - y_{\mathrm{obj}}\|^2$, $(\lambda_g, \lambda_k, \lambda_o) = (0.1, 0.05, 0.05)$. 선택적으로 불확실성 가중 \cite{kendall2018multi}을 지원한다.
  \item[T7 라벨 평활화 \cite{szegedy2016rethinking}] $y_{\mathrm{smooth}} = y(1-\epsilon) + \epsilon/2$, $\epsilon = 0.05$. 균등 분포 방향의 KL 정규화로 과신을 줄인다 \cite{muller2019when}.
\end{description}

이 밖에 보간·검출기 절제로 T8(보간 없이 60 초 기준선), T9(킬 궤적 덮어쓰기 없는 검출기), T10(60 초 격자만 사용하는 검출기)을 등록하여 제 \ref{ch:localization} 장의 설계 요소를 같은 프로토콜로 평가할 수 있게 하였다.
